% Options for packages loaded elsewhere
% Options for packages loaded elsewhere
\PassOptionsToPackage{unicode}{hyperref}
\PassOptionsToPackage{hyphens}{url}
\PassOptionsToPackage{dvipsnames,svgnames,x11names}{xcolor}
%
\documentclass[
  11pt,
  a4paper,
  DIV=11,
  numbers=noendperiod]{scrartcl}
\usepackage{xcolor}
\usepackage{amsmath,amssymb}
\setcounter{secnumdepth}{-\maxdimen} % remove section numbering
\usepackage{iftex}
\ifPDFTeX
  \usepackage[T1]{fontenc}
  \usepackage[utf8]{inputenc}
  \usepackage{textcomp} % provide euro and other symbols
\else % if luatex or xetex
  \usepackage{unicode-math} % this also loads fontspec
  \defaultfontfeatures{Scale=MatchLowercase}
  \defaultfontfeatures[\rmfamily]{Ligatures=TeX,Scale=1}
\fi
\usepackage{lmodern}
\ifPDFTeX\else
  % xetex/luatex font selection
  \setmainfont[]{Avenir Next}
\fi
% Use upquote if available, for straight quotes in verbatim environments
\IfFileExists{upquote.sty}{\usepackage{upquote}}{}
\IfFileExists{microtype.sty}{% use microtype if available
  \usepackage[]{microtype}
  \UseMicrotypeSet[protrusion]{basicmath} % disable protrusion for tt fonts
}{}
\makeatletter
\@ifundefined{KOMAClassName}{% if non-KOMA class
  \IfFileExists{parskip.sty}{%
    \usepackage{parskip}
  }{% else
    \setlength{\parindent}{0pt}
    \setlength{\parskip}{6pt plus 2pt minus 1pt}}
}{% if KOMA class
  \KOMAoptions{parskip=half}}
\makeatother
% Make \paragraph and \subparagraph free-standing
\makeatletter
\ifx\paragraph\undefined\else
  \let\oldparagraph\paragraph
  \renewcommand{\paragraph}{
    \@ifstar
      \xxxParagraphStar
      \xxxParagraphNoStar
  }
  \newcommand{\xxxParagraphStar}[1]{\oldparagraph*{#1}\mbox{}}
  \newcommand{\xxxParagraphNoStar}[1]{\oldparagraph{#1}\mbox{}}
\fi
\ifx\subparagraph\undefined\else
  \let\oldsubparagraph\subparagraph
  \renewcommand{\subparagraph}{
    \@ifstar
      \xxxSubParagraphStar
      \xxxSubParagraphNoStar
  }
  \newcommand{\xxxSubParagraphStar}[1]{\oldsubparagraph*{#1}\mbox{}}
  \newcommand{\xxxSubParagraphNoStar}[1]{\oldsubparagraph{#1}\mbox{}}
\fi
\makeatother

\usepackage{color}
\usepackage{fancyvrb}
\newcommand{\VerbBar}{|}
\newcommand{\VERB}{\Verb[commandchars=\\\{\}]}
\DefineVerbatimEnvironment{Highlighting}{Verbatim}{commandchars=\\\{\}}
% Add ',fontsize=\small' for more characters per line
\usepackage{framed}
\definecolor{shadecolor}{RGB}{241,243,245}
\newenvironment{Shaded}{\begin{snugshade}}{\end{snugshade}}
\newcommand{\AlertTok}[1]{\textcolor[rgb]{0.68,0.00,0.00}{#1}}
\newcommand{\AnnotationTok}[1]{\textcolor[rgb]{0.37,0.37,0.37}{#1}}
\newcommand{\AttributeTok}[1]{\textcolor[rgb]{0.40,0.45,0.13}{#1}}
\newcommand{\BaseNTok}[1]{\textcolor[rgb]{0.68,0.00,0.00}{#1}}
\newcommand{\BuiltInTok}[1]{\textcolor[rgb]{0.00,0.23,0.31}{#1}}
\newcommand{\CharTok}[1]{\textcolor[rgb]{0.13,0.47,0.30}{#1}}
\newcommand{\CommentTok}[1]{\textcolor[rgb]{0.37,0.37,0.37}{#1}}
\newcommand{\CommentVarTok}[1]{\textcolor[rgb]{0.37,0.37,0.37}{\textit{#1}}}
\newcommand{\ConstantTok}[1]{\textcolor[rgb]{0.56,0.35,0.01}{#1}}
\newcommand{\ControlFlowTok}[1]{\textcolor[rgb]{0.00,0.23,0.31}{\textbf{#1}}}
\newcommand{\DataTypeTok}[1]{\textcolor[rgb]{0.68,0.00,0.00}{#1}}
\newcommand{\DecValTok}[1]{\textcolor[rgb]{0.68,0.00,0.00}{#1}}
\newcommand{\DocumentationTok}[1]{\textcolor[rgb]{0.37,0.37,0.37}{\textit{#1}}}
\newcommand{\ErrorTok}[1]{\textcolor[rgb]{0.68,0.00,0.00}{#1}}
\newcommand{\ExtensionTok}[1]{\textcolor[rgb]{0.00,0.23,0.31}{#1}}
\newcommand{\FloatTok}[1]{\textcolor[rgb]{0.68,0.00,0.00}{#1}}
\newcommand{\FunctionTok}[1]{\textcolor[rgb]{0.28,0.35,0.67}{#1}}
\newcommand{\ImportTok}[1]{\textcolor[rgb]{0.00,0.46,0.62}{#1}}
\newcommand{\InformationTok}[1]{\textcolor[rgb]{0.37,0.37,0.37}{#1}}
\newcommand{\KeywordTok}[1]{\textcolor[rgb]{0.00,0.23,0.31}{\textbf{#1}}}
\newcommand{\NormalTok}[1]{\textcolor[rgb]{0.00,0.23,0.31}{#1}}
\newcommand{\OperatorTok}[1]{\textcolor[rgb]{0.37,0.37,0.37}{#1}}
\newcommand{\OtherTok}[1]{\textcolor[rgb]{0.00,0.23,0.31}{#1}}
\newcommand{\PreprocessorTok}[1]{\textcolor[rgb]{0.68,0.00,0.00}{#1}}
\newcommand{\RegionMarkerTok}[1]{\textcolor[rgb]{0.00,0.23,0.31}{#1}}
\newcommand{\SpecialCharTok}[1]{\textcolor[rgb]{0.37,0.37,0.37}{#1}}
\newcommand{\SpecialStringTok}[1]{\textcolor[rgb]{0.13,0.47,0.30}{#1}}
\newcommand{\StringTok}[1]{\textcolor[rgb]{0.13,0.47,0.30}{#1}}
\newcommand{\VariableTok}[1]{\textcolor[rgb]{0.07,0.07,0.07}{#1}}
\newcommand{\VerbatimStringTok}[1]{\textcolor[rgb]{0.13,0.47,0.30}{#1}}
\newcommand{\WarningTok}[1]{\textcolor[rgb]{0.37,0.37,0.37}{\textit{#1}}}

\usepackage{longtable,booktabs,array}
\usepackage{calc} % for calculating minipage widths
% Correct order of tables after \paragraph or \subparagraph
\usepackage{etoolbox}
\makeatletter
\patchcmd\longtable{\par}{\if@noskipsec\mbox{}\fi\par}{}{}
\makeatother
% Allow footnotes in longtable head/foot
\IfFileExists{footnotehyper.sty}{\usepackage{footnotehyper}}{\usepackage{footnote}}
\makesavenoteenv{longtable}
\usepackage{graphicx}
\makeatletter
\newsavebox\pandoc@box
\newcommand*\pandocbounded[1]{% scales image to fit in text height/width
  \sbox\pandoc@box{#1}%
  \Gscale@div\@tempa{\textheight}{\dimexpr\ht\pandoc@box+\dp\pandoc@box\relax}%
  \Gscale@div\@tempb{\linewidth}{\wd\pandoc@box}%
  \ifdim\@tempb\p@<\@tempa\p@\let\@tempa\@tempb\fi% select the smaller of both
  \ifdim\@tempa\p@<\p@\scalebox{\@tempa}{\usebox\pandoc@box}%
  \else\usebox{\pandoc@box}%
  \fi%
}
% Set default figure placement to htbp
\def\fps@figure{htbp}
\makeatother





\setlength{\emergencystretch}{3em} % prevent overfull lines

\providecommand{\tightlist}{%
  \setlength{\itemsep}{0pt}\setlength{\parskip}{0pt}}



 


\usepackage[document]{ragged2e}
\KOMAoption{captions}{tableheading}
\makeatletter
\@ifpackageloaded{caption}{}{\usepackage{caption}}
\AtBeginDocument{%
\ifdefined\contentsname
  \renewcommand*\contentsname{Table of contents}
\else
  \newcommand\contentsname{Table of contents}
\fi
\ifdefined\listfigurename
  \renewcommand*\listfigurename{List of Figures}
\else
  \newcommand\listfigurename{List of Figures}
\fi
\ifdefined\listtablename
  \renewcommand*\listtablename{List of Tables}
\else
  \newcommand\listtablename{List of Tables}
\fi
\ifdefined\figurename
  \renewcommand*\figurename{Figure}
\else
  \newcommand\figurename{Figure}
\fi
\ifdefined\tablename
  \renewcommand*\tablename{Table}
\else
  \newcommand\tablename{Table}
\fi
}
\@ifpackageloaded{float}{}{\usepackage{float}}
\floatstyle{ruled}
\@ifundefined{c@chapter}{\newfloat{codelisting}{h}{lop}}{\newfloat{codelisting}{h}{lop}[chapter]}
\floatname{codelisting}{Listing}
\newcommand*\listoflistings{\listof{codelisting}{List of Listings}}
\makeatother
\makeatletter
\makeatother
\makeatletter
\@ifpackageloaded{caption}{}{\usepackage{caption}}
\@ifpackageloaded{subcaption}{}{\usepackage{subcaption}}
\makeatother
\usepackage{bookmark}
\IfFileExists{xurl.sty}{\usepackage{xurl}}{} % add URL line breaks if available
\urlstyle{same}
\hypersetup{
  pdftitle={12. Iterationen},
  colorlinks=true,
  linkcolor={blue},
  filecolor={Maroon},
  citecolor={Blue},
  urlcolor={Blue},
  pdfcreator={LaTeX via pandoc}}


\title{12. Iterationen}
\author{}
\date{}
\begin{document}
\maketitle


\textbf{Problem:} Eine bestimmte Anzahl von Primzahlen (Liste von
Primzahlen) soll multipliziert werden. Die Primzahlen sind in einer also
in einer Liste gespeichert.

\begin{Shaded}
\begin{Highlighting}[numbers=left,,]
\NormalTok{primzahlen }\OperatorTok{=}\NormalTok{ [}\DecValTok{2}\NormalTok{,}\DecValTok{3}\NormalTok{,}\DecValTok{5}\NormalTok{,}\DecValTok{7}\NormalTok{,}\DecValTok{11}\NormalTok{,}\DecValTok{13}\NormalTok{]}

\NormalTok{product }\OperatorTok{=}\NormalTok{ primzahlen[}\DecValTok{0}\NormalTok{]}\OperatorTok{*}\NormalTok{ primzahlen[}\DecValTok{1}\NormalTok{]}\OperatorTok{*}\NormalTok{primzahlen[}\DecValTok{2}\NormalTok{]}\OperatorTok{*}\NormalTok{primzahlen[}\DecValTok{3}\NormalTok{]}\OperatorTok{*}\NormalTok{primzahlen[}\DecValTok{4}\NormalTok{]}\OperatorTok{*}\NormalTok{primzahlen[}\DecValTok{5}\NormalTok{]. }
\BuiltInTok{print}\NormalTok{(product) }
\end{Highlighting}
\end{Shaded}

\begin{verbatim}
30030
\end{verbatim}

\$\Rightarrow \$ Dieser Code ist schreibaufwendig und unübersichtlich\\
\$\Rightarrow \$ Umständlich, bei großen Listen

andere Lösung: Interationen

\subsubsection{\texorpdfstring{12.1
\texttt{for}-Schleife}{12.1 for-Schleife}}\label{for-schleife}

\begin{Shaded}
\begin{Highlighting}[numbers=left,,]
\NormalTok{primzahlen }\OperatorTok{=}\NormalTok{ [}\DecValTok{2}\NormalTok{,}\DecValTok{3}\NormalTok{,}\DecValTok{5}\NormalTok{,}\DecValTok{7}\NormalTok{,}\DecValTok{11}\NormalTok{,}\DecValTok{13}\NormalTok{]}
\NormalTok{product }\OperatorTok{=} \DecValTok{1}
\ControlFlowTok{for}\NormalTok{ number }\KeywordTok{in}\NormalTok{ primzahlen:}
\NormalTok{    product }\OperatorTok{=}\NormalTok{ product }\OperatorTok{*}\NormalTok{ number }\CommentTok{\#same as product *= number}

\BuiltInTok{print}\NormalTok{ (product)}
\end{Highlighting}
\end{Shaded}

\begin{verbatim}
30030
\end{verbatim}

\textbf{Syntax}

\begin{Shaded}
\begin{Highlighting}[numbers=left,,]
\ControlFlowTok{for}\NormalTok{ var }\KeywordTok{in}\NormalTok{ expresssion:   }
\NormalTok{        block}
\end{Highlighting}
\end{Shaded}

\begin{itemize}
\tightlist
\item
  Schlüsselworte sind \texttt{for}und \texttt{in}
\item
  \begin{enumerate}
  \def\labelenumi{\arabic{enumi}.}
  \tightlist
  \item
    Zeile heißt Schleifenkopf
  \end{enumerate}
\item
  \begin{enumerate}
  \def\labelenumi{\arabic{enumi}.}
  \setcounter{enumi}{1}
  \tightlist
  \item
    Zeile Schleifenrumpf \texttt{block} mit einer oder mehreren
    Anweisungen
  \end{enumerate}
\item
  Schleifenvariable \texttt{var} im Schleifenkopf
\item
  Schleifeniteration ist ein Durchlauf (Ausführnung eines
  Schleifenrumpfs)
\end{itemize}

\textbf{Beispiel:}

\begin{Shaded}
\begin{Highlighting}[numbers=left,,]
\ControlFlowTok{for}\NormalTok{ j }\KeywordTok{in} \StringTok{"Hallo"}\NormalTok{:}
    \BuiltInTok{print}\NormalTok{(j}\OperatorTok{*} \DecValTok{2}\NormalTok{)}
\end{Highlighting}
\end{Shaded}

\begin{verbatim}
HH
aa
ll
ll
oo
\end{verbatim}

\begin{Shaded}
\begin{Highlighting}[numbers=left,,]
\ControlFlowTok{for}\NormalTok{ ingredient }\KeywordTok{in}\NormalTok{ (}\StringTok{"sugar"}\NormalTok{, }\StringTok{"water"}\NormalTok{, }\StringTok{"oil"}\NormalTok{):}
    \ControlFlowTok{if}\NormalTok{ ingredient }\OperatorTok{==} \StringTok{"sugar"}\NormalTok{:}
        \BuiltInTok{print}\NormalTok{(}\StringTok{"sweet!"}\NormalTok{)}
\end{Highlighting}
\end{Shaded}

\begin{verbatim}
sweet!
\end{verbatim}

\textbf{Ablauf der Schleife beeinflussen}

\begin{itemize}
\tightlist
\item
  \texttt{break} - Schleifenrumpf wird umehend beendet (eventuell
  vorzeitg)
\item
  \texttt{continue} - die aktuelle Schleifeniteration wird vorzeitig
  beendet, es wird in den Schleifenkopf gesprungen und die
  Schleifenvariable wird auf den nächsten Wert gesetzt.
\item
  \texttt{else-}Zweig - wird nach Beendigungder Schleife ausgeführt,
  genau dann wenn die Schleife nicht mit \texttt{break} verlassen wurde.
\end{itemize}

\textbf{Beispiel}

\begin{Shaded}
\begin{Highlighting}[numbers=left,,]
\NormalTok{noten\_und\_anzahl }\OperatorTok{=}\NormalTok{[(}\StringTok{"sehr gut"}\NormalTok{, }\DecValTok{1}\NormalTok{), (}\StringTok{"gut"}\NormalTok{, }\DecValTok{3}\NormalTok{), (}\StringTok{"befriedigend"}\NormalTok{,}\DecValTok{0}\NormalTok{), (}\StringTok{"ausreichend"}\NormalTok{, }\DecValTok{2}\NormalTok{)]}

\ControlFlowTok{for}\NormalTok{ na }\KeywordTok{in}\NormalTok{ noten\_und\_anzahl:}
\NormalTok{    note , anzahl }\OperatorTok{=}\NormalTok{ na}
    \ControlFlowTok{if}\NormalTok{ anzahl }\OperatorTok{==} \DecValTok{0}\NormalTok{:}
        \ControlFlowTok{continue}
    \ControlFlowTok{if}\NormalTok{ note }\OperatorTok{==} \StringTok{"h"}\NormalTok{:}
        \BuiltInTok{print}\NormalTok{(}\StringTok{"Tolle Leistung"}\NormalTok{)}
        \ControlFlowTok{break}
\ControlFlowTok{else}\NormalTok{:}
    \BuiltInTok{print}\NormalTok{(}\StringTok{"Durchschnitt war in Ordnung"}\NormalTok{)}
     
\end{Highlighting}
\end{Shaded}

\begin{verbatim}
Durchschnitt war in Ordnung
\end{verbatim}

Alternativ kann man das auch so aufschreiben:

\begin{Shaded}
\begin{Highlighting}[numbers=left,,]
\NormalTok{noten\_und\_anzahl }\OperatorTok{=}\NormalTok{[(}\StringTok{"sehr gut"}\NormalTok{, }\DecValTok{1}\NormalTok{), (}\StringTok{"gut"}\NormalTok{, }\DecValTok{3}\NormalTok{), (}\StringTok{"befriedigend"}\NormalTok{,}\DecValTok{0}\NormalTok{), (}\StringTok{"ausreichend, 2"}\NormalTok{)]}

\ControlFlowTok{for}\NormalTok{ note, anzahl }\KeywordTok{in}\NormalTok{ noten\_und\_anzahl:}
    \ControlFlowTok{if}\NormalTok{ anzahl }\OperatorTok{==}\DecValTok{0}\NormalTok{:}
        \ControlFlowTok{continue}
    \ControlFlowTok{if}\NormalTok{ note }\OperatorTok{==} \StringTok{"gut"}\NormalTok{:}
        \BuiltInTok{print}\NormalTok{(}\StringTok{"Tolle Leistung"}\NormalTok{)}
        \ControlFlowTok{break}
\ControlFlowTok{else}\NormalTok{:}
    \BuiltInTok{print}\NormalTok{(}\StringTok{"Durchschnitt war in Ordnung"}\NormalTok{)}
     
\end{Highlighting}
\end{Shaded}

\begin{verbatim}
Tolle Leistung
\end{verbatim}

\textbf{Nützliche Funktionen}

\begin{itemize}
\tightlist
\item
  \texttt{range}
\item
  \texttt{zip}
\item
  \texttt{reversed}
\end{itemize}

\subsubsection{\texorpdfstring{12.2
\texttt{range}}{12.2 range}}\label{range}

\begin{itemize}
\tightlist
\item
  erzeugt eine Folge von Indices für die Schleifendurchläufe, genauer
  einen Iterator.
\item
  \texttt{range(stop)} - ergibt 0, 1, \ldots, stop -1
\item
  \texttt{range(start,\ stop)} - ergibt start, start + 1, \ldots{} ,
  stop -1
\item
  \texttt{range(start,\ stop,\ step)} - ergibt start, start +step,
  start+ 2* step, \ldots{} , stop -1
\end{itemize}

\textbf{Beispiele:}

\begin{Shaded}
\begin{Highlighting}[numbers=left,,]
\BuiltInTok{range}\NormalTok{(}\DecValTok{5}\NormalTok{)}
\end{Highlighting}
\end{Shaded}

\begin{verbatim}
range(0, 5)
\end{verbatim}

\begin{Shaded}
\begin{Highlighting}[numbers=left,,]
\BuiltInTok{print}\NormalTok{(}\BuiltInTok{range}\NormalTok{(}\DecValTok{2}\NormalTok{, }\DecValTok{30}\NormalTok{, }\DecValTok{5}\NormalTok{)) }\CommentTok{\#keine Liste, print gibt keine Ausgabe}
\end{Highlighting}
\end{Shaded}

\begin{verbatim}
range(2, 30, 5)
\end{verbatim}

\begin{Shaded}
\begin{Highlighting}[numbers=left,,]
\BuiltInTok{list}\NormalTok{(}\BuiltInTok{range}\NormalTok{(}\DecValTok{0}\NormalTok{,}\DecValTok{9}\NormalTok{))}
\end{Highlighting}
\end{Shaded}

\begin{verbatim}
[0, 1, 2, 3, 4, 5, 6, 7, 8]
\end{verbatim}

\begin{Shaded}
\begin{Highlighting}[numbers=left,,]
\BuiltInTok{list}\NormalTok{ (}\BuiltInTok{range}\NormalTok{(}\DecValTok{2}\NormalTok{, }\DecValTok{30}\NormalTok{, }\DecValTok{5}\NormalTok{))}
\end{Highlighting}
\end{Shaded}

\begin{verbatim}
[2, 7, 12, 17, 22, 27]
\end{verbatim}

\begin{Shaded}
\begin{Highlighting}[numbers=left,,]
\ControlFlowTok{for}\NormalTok{ i }\KeywordTok{in} \BuiltInTok{range}\NormalTok{(}\DecValTok{3}\NormalTok{,}\DecValTok{6}\NormalTok{):}
    \BuiltInTok{print}\NormalTok{ (i, }\StringTok{"**3 = "}\NormalTok{, i }\OperatorTok{**} \DecValTok{3}\NormalTok{)}
\end{Highlighting}
\end{Shaded}

\begin{verbatim}
3 **3 =  27
4 **3 =  64
5 **3 =  125
\end{verbatim}

\subsubsection{\texorpdfstring{12.3 \texttt{zip}}{12.3 zip}}\label{zip}

\begin{itemize}
\tightlist
\item
  Funktion
\item
  nimmt eine oder mehrere Sequenzen und liefert eine ``Liste'' von
  Tupeln mit korrespndierenden Elementen
\item
  erzeugt eigentlich keine Liste, sondern einen Iterator
\end{itemize}

\textbf{Beipsiel}

\begin{Shaded}
\begin{Highlighting}[numbers=left,,]
\NormalTok{fleisch }\OperatorTok{=}\NormalTok{ [}\StringTok{"Rind"}\NormalTok{, }\StringTok{"Huhn"}\NormalTok{, }\StringTok{"Hund"}\NormalTok{]}
\NormalTok{beilage }\OperatorTok{=}\NormalTok{[}\StringTok{"Kartoffeln"}\NormalTok{, }\StringTok{"Pommes"}\NormalTok{, }\StringTok{"Spätzle"}\NormalTok{]}

\BuiltInTok{print}\NormalTok{(}\BuiltInTok{list}\NormalTok{(}\BuiltInTok{zip}\NormalTok{(fleisch, beilage)))}
\end{Highlighting}
\end{Shaded}

\begin{verbatim}
[('Rind', 'Kartoffeln'), ('Huhn', 'Pommes'), ('Hund', 'Spätzle')]
\end{verbatim}

\begin{itemize}
\tightlist
\item
  mit \texttt{zip} können mehrere Sequenzen parallel durchlaufen werden
\item
  bei unterschiedlicher Länge ist das Ergbnis so lang, wie die kürzere
  Sequenz.
\end{itemize}

\textbf{Beispiel}

\begin{Shaded}
\begin{Highlighting}[numbers=left,,]
\ControlFlowTok{for}\NormalTok{ xyz }\KeywordTok{in} \BuiltInTok{zip}\NormalTok{(}\StringTok{"rind"}\NormalTok{, }\StringTok{"huhn"}\NormalTok{, }\BuiltInTok{range}\NormalTok{(}\DecValTok{5}\NormalTok{,}\DecValTok{10}\NormalTok{)):}
\NormalTok{    x, y, z }\OperatorTok{=}\NormalTok{ xyz}
    \BuiltInTok{print}\NormalTok{(x, y, z)}
\end{Highlighting}
\end{Shaded}

\begin{verbatim}
r h 5
i u 6
n h 7
d n 8
\end{verbatim}

\subsubsection{\texorpdfstring{12.4
\texttt{reversed}}{12.4 reversed}}\label{reversed}

\begin{itemize}
\tightlist
\item
  Funktion
\item
  Durchlauf einer Sequenz in umgekehrter Richtung
\end{itemize}

\textbf{Beispiel}

\begin{Shaded}
\begin{Highlighting}[numbers=left,,]
\ControlFlowTok{for}\NormalTok{ x }\KeywordTok{in} \BuiltInTok{reversed}\NormalTok{(}\StringTok{"hallo"}\NormalTok{):}
    \BuiltInTok{print}\NormalTok{(x)}
\end{Highlighting}
\end{Shaded}

\begin{verbatim}
o
l
l
a
h
\end{verbatim}

\subsubsection{12.5 Beispiel: Implementierung einer
Fakultätsfunktion}\label{beispiel-implementierung-einer-fakultuxe4tsfunktion}

Fakultät ist wie folgt definiert:

\[0! = 1\]

\[(n+1)! = (n+1)\cdot n!\]

Aufgabe: Entwickle eine Funktion \texttt{fact}, welche die Fakulät einer
ganzen Zahl berechnet.

\begin{enumerate}
\def\labelenumi{\arabic{enumi}.}
\item
  Bezeichner und Datentypen

  \begin{itemize}
  \tightlist
  \item
    Die Eingabe ist \texttt{n:\ int} (mit n \textgreater= 0)
  \item
    Die Ausgabe ist \texttt{int}
  \end{itemize}
\item
  Funktionsgerüst
  \texttt{python\ \ \ \ \ \ \ \ \ def\ fact(n:\ int)-\textgreater{}\ int\ \#\#assume\ n\textgreater{}1\ \ \ \ \ \ \ \ \ \#\ fill\ in\ \ \ \ \ \ \ \ \ return}
\item
  Beispiele

  \texttt{assert\ fact(0)\ ==\ 1}\strut \\
  \texttt{assert\ fact(1)\ ==\ 1}\strut \\
  \texttt{assert\ fact(3)\ ==\ 6}
\item
  Ergebnis
\end{enumerate}

\begin{Shaded}
\begin{Highlighting}[numbers=left,,]
\KeywordTok{def}\NormalTok{ fact(}
\NormalTok{        n: }\BuiltInTok{int} 
\NormalTok{        ) }\OperatorTok{{-}\textgreater{}} \BuiltInTok{int}\NormalTok{:}
\NormalTok{    result }\OperatorTok{=} \DecValTok{1}
    \ControlFlowTok{for}\NormalTok{ i }\KeywordTok{in} \BuiltInTok{range}\NormalTok{ (}\DecValTok{1}\NormalTok{, n}\OperatorTok{+}\DecValTok{1}\NormalTok{):}
\NormalTok{        result}\OperatorTok{=}\NormalTok{ result }\OperatorTok{*}\NormalTok{ i}
    \ControlFlowTok{return}\NormalTok{ result}
\end{Highlighting}
\end{Shaded}

\begin{Shaded}
\begin{Highlighting}[numbers=left,,]
\BuiltInTok{print}\NormalTok{(fact(}\DecValTok{10}\NormalTok{))}
\end{Highlighting}
\end{Shaded}

\begin{verbatim}
3628800
\end{verbatim}




\end{document}
